\section{Proofs for ranked models and finite rational closure}
\label{annex:p08-ranked-models-rational-closure}

\begin{annexitem}{annex:p08.rational.modularity}{thm:p08-rational-modularity}{Quotient ranking and the exact smoothness condition}
\begin{proof}
Number the four conditions in their stated order. Suppose condition 1
holds and \(x\prec y\), \(z\not\prec y\). If \(y\prec z\),
transitivity gives \(x\prec z\). Otherwise \(y,z\) are
incomparable, and their common strict predecessors include \(x\).
Thus \(x\prec z\), proving modularity. Modularity gives the
contrapositive of condition 3: if \(x\prec z\), then
\(y\prec z\) or \(x\prec y\).

Under condition 3, write \(x\approx y\) for incomparability,
including equality. It is reflexive and symmetric. If
\(x\approx y\approx z\), condition 3 applied in both directions
gives \(x\not\prec z\) and \(z\not\prec x\), so it is
transitive. Condition 3 also implies modularity by the same
contrapositive. If \(x\prec y\), \(x'\approx x\) and
\(y'\approx y\), modularity first gives \(x'\prec y\), since
\(x\prec x'\) is impossible, and then \(x'\prec y'\), since
\(y'\prec y\) is impossible. The order between incomparable
classes is consequently well defined. Any two different classes are
comparable; transitivity and irreflexivity descend to the quotient.
The map sending a point to its class induces the original preference,
proving condition 4. Conversely a ranking makes incomparable elements
have equal ranks, hence identical strict predecessors, proving 1.

For a ranked set \(T\), a point \(t\in T\) is minimal iff
there is no \(s\in T\) of strictly smaller rank. Total ordering
makes this equivalent to \(r(t)\) being the least element of
\(r[T]\). If that least rank exists, every point not at it lies
above every point at it, so \(T\) is smooth. If a nonempty \(T\)
is smooth, applying smoothness to any point supplies a minimal point,
whose rank is least. These statements apply individually to each
formula-definable set; no arbitrary-subset or well-order assumption
has been introduced.
\end{proof}
\end{annexitem}

\begin{annexitem}{annex:p08.rational.soundness}{thm:p08-rational-soundness}{The ranked negative-premise argument}
\begin{proof}
Theorem~\ref{thm:p08-klm-soundness}, with the full classical valuation
universe, proves all P rules. Suppose
\(\alpha\mathrel{|\!\sim}_W\gamma\) and
\(\alpha\mathrel{|\!\not\sim}_W\neg\beta\).
The second statement supplies \(t\) minimal in
\(\widehat\alpha\) with \(t\equiv\beta\).
For any \(s\) minimal in \(\widehat{\alpha\wedge\beta}\),
minimality and the presence of \(t\) give \(r(s)\leq r(t)\).
If an \(\alpha\)-state \(u\) preceded \(s\), its rank
would be strictly smaller than \(r(t)\), contradicting minimality
of \(t\) in \(\widehat\alpha\). Hence \(s\) is minimal
in \(\widehat\alpha\) and satisfies \(\gamma\).
This establishes Rational Monotonicity. The argument itself does not
need to infer existence of minima from smoothness: \(t\) is already
given by the negative premise and the conclusion quantifies over
whatever minima exist. The P soundness used earlier does require
smoothness. If \(\widehat\alpha\) is empty, all assertions with
that antecedent hold, so the negative premise cannot hold.
\end{proof}
\end{annexitem}

\begin{annexitem}{annex:p08.rational.transport}{thm:p08-rational-transport}{Normal witnesses, rational comparison and every quotient dependency}
\begin{proof}
First supply the normal-world argument. If a normal world exists,
\(\alpha\mathrel{|\!\sim}\bot\) cannot hold, since no world
satisfies \(\bot\). If \(\alpha\mathrel{|\!\not\sim}\delta\),
let \(T\) consist of every consequent of \(\alpha\), together
with \(\neg\delta\). If it were finitely unsatisfiable, finitely
many consequents would entail \(\delta\). And and Right Weakening
would give \(\alpha\mathrel{|\!\sim}\delta\), a contradiction;
the empty-conjunction case follows from Reflexivity and Right Weakening
to \(\top\). Classical compactness therefore supplies a normal
world falsifying \(\delta\). In particular \(\delta=\bot\)
proves existence for every consistent formula. This also establishes
the witness for a missing assertion without assuming finiteness of
the consequence theory.

The four derived P rules in the sheet can be verified using the already
proved preferential representation, without using ranked representation.
In a smooth model, \(\alpha\mathrel{|\!\sim}\bot\) means
\(\widehat\alpha\) is empty, proving rule (14). For rule (15),
suppose every minimum of \(\widehat{\alpha\vee\beta}\)
falsifies \(\beta\). If a minimal \(\alpha\)-state were not
minimal in the union, smoothness would supply a union-minimum below
it. That minimum falsifies \(\beta\), hence satisfies \(\alpha\),
a contradiction. Thus every minimal \(\alpha\)-state falsifies
\(\beta\). For rule (16), a minimum of
\(\widehat{\beta\vee\gamma}\) that is not a minimum of the
larger triple union would have a triple-union minimum below it.
That predecessor falsifies both \(\alpha,\beta\), hence satisfies
\(\gamma\), again a contradiction. Triple-union minima already
falsify \(\beta\), proving the result. For rule (17), a
triple-union minimum satisfying \(\alpha\) would also be a
minimum of \(\widehat{\alpha\vee\beta}\), contradicting its
premise. These arguments include empty sets and prove each displayed
derived rule for every P relation by exact representation.

Suppose \(\alpha\vee\gamma\mathrel{|\!\sim}\neg\alpha\)
and \(\beta\vee\gamma\mathrel{|\!\not\sim}\neg\beta\).
Rule (17) gives \(D=\alpha\vee\beta\vee\gamma\mathrel{|\!\sim}\neg\alpha\).
If \(D\mathrel{|\!\sim}\neg(\alpha\vee\beta)\), Right
Weakening and rule (16) would give the forbidden
\(\beta\vee\gamma\mathrel{|\!\sim}\neg\beta\).
Thus \(D\mathrel{|\!\not\sim}\neg(\alpha\vee\beta)\).
Rational Monotonicity restricts \(D\) to \(\alpha\vee\beta\)
and gives \(\alpha\vee\beta\mathrel{|\!\sim}\neg\alpha\).
This is the source's comparison inference (18).

If \(\alpha\mathcal R\beta\) and \(\beta\mathcal R\gamma\)
but not \(\alpha\mathcal R\gamma\), inference (18) gives
\(\alpha\vee\beta\mathrel{|\!\sim}\neg\alpha\),
contradicting the first assumption. Hence \(\mathcal R\) is
transitive. If neither direction held between \(\alpha,\beta\in S\),
And with Reflexivity would give
\(\alpha\vee\beta\mathrel{|\!\sim}\bot\).
Rule (14), restricting to \(\beta\), would contradict \(\beta\in S\).
This proves totality and, setting the formulas equal, reflexivity.
Mutual comparison is consequently an equivalence relation. Transitivity
shows that changing either representative preserves comparison, so
the quotient order is well defined, antisymmetric and total.

For normality transport, let \(\beta\mathrel{|\!\sim}\gamma\).
The derived S rule gives
\(\alpha\vee\beta\mathrel{|\!\sim}\beta\to\gamma\).
Since \(\alpha\mathcal R\beta\), Rational Monotonicity and
Left Logical Equivalence give
\(\alpha\mathrel{|\!\sim}\beta\to\gamma\).
A normal \(\alpha\)-world satisfying \(\beta\) thus satisfies
every such \(\gamma\) and is normal for \(\beta\).
If \(\overline\beta<\overline\alpha\), failure of
\(\alpha\mathcal R\beta\) is
\(\alpha\vee\beta\mathrel{|\!\sim}\neg\alpha\);
rule (15) gives \(\beta\mathrel{|\!\sim}\neg\alpha\).

Finally, a normal \(\alpha\)-world satisfying \(\beta\)
forces \(\beta\) to be consistent. Otherwise rule (14) gives
\(\alpha\wedge\beta\mathrel{|\!\sim}\bot\), and S and
Right Weakening give \(\alpha\mathrel{|\!\sim}\neg\beta\),
contradicting that world. If
\(\overline\alpha<\overline\beta\), the previous paragraph
likewise gives \(\alpha\mathrel{|\!\sim}\neg\beta\).
Totality therefore forces
\(\overline\beta\leq\overline\alpha\), as required.
\end{proof}
\end{annexitem}

\begin{annexitem}{annex:p08.rational.minima}{thm:p08-rational-canonical-minima}{All definable minima, and the empty inconsistent case}
\begin{proof}
If \(\langle m,\alpha\rangle\) satisfies \(\beta\), the
last transport conclusion implies \(\beta\in S\) and
\(\overline\beta\leq\overline\alpha\). In particular an
inconsistent \(\beta\) has no states in its definable set.
For consistent \(\beta\), choose a normal world \(n\) for it.
The state \(\langle n,\beta\rangle\) exists and satisfies
\(\beta\). Minimality of \(\langle m,\alpha\rangle\)
in \(\widehat\beta\) rules out
\(\overline\beta<\overline\alpha\). Together with the previous
inequality this proves equality of classes.

Conversely suppose \(m\models\beta\) and the classes are equal.
Any state \(\langle n,\gamma\rangle\) preceding it has
\(\overline\gamma<\overline\beta\). The strict-class conclusion
gives \(\gamma\mathrel{|\!\sim}\neg\beta\), so its normal
world \(n\) cannot satisfy \(\beta\). There is therefore no
predecessor in \(\widehat\beta\), proving minimality.

For smoothness take any state in \(\widehat\beta\). Its class
is at least \(\overline\beta\). Equality gives minimality;
strict inequality gives the minimal predecessor
\(\langle n,\beta\rangle\) just constructed. Empty definable
sets are smooth vacuously. This proves smoothness for all formulas,
including inconsistent ones whose class was never defined.
\end{proof}
\end{annexitem}

\begin{annexitem}{annex:p08.rational.representation}{thm:p08-rational-general-representation}{Preservation, reflection and the finite logical quotient}
\begin{proof}
Ranked soundness proves one direction. For the other, the set \(S\)
of consistent formulas, set \(E\) of equivalence classes, and every
normal-world/formula pair are genuine sets. The class order is total
by item~\ref{annex:p08.rational.transport}; its strict part therefore
gives a strict partial order on states via the rank map. The preceding
item proves the required smoothness.

If \(\alpha\mathrel{|\!\sim}\beta\), consider a minimal
\(\alpha\)-state \(\langle m,\gamma\rangle\). Its class
equals \(\overline\alpha\), hence
\(\gamma\mathcal R\alpha\). Normality transport makes \(m\)
normal for \(\alpha\), and so it satisfies \(\beta\).
If there are no \(\alpha\)-states, the assertion holds vacuously.
If instead \(\alpha\mathrel{|\!\not\sim}\beta\), the
compactness witness is a normal world \(m\) for \(\alpha\)
falsifying \(\beta\). It forces \(\alpha\) to be consistent,
and the state \(\langle m,\alpha\rangle\) is minimal for
\(\alpha\) by its equal class. It is the required counterstate.
Thus the relation is preserved and reflected on every formula pair,
without choosing one normal world per formula or assuming ordinal ranks.

For \(d\) atoms there are \(2^d\) valuations and at most
\(2^{2^d}\) classical formula-equivalence classes. Choose one formula
per such class and retain only these formula coordinates among consistent
ones. Left Logical Equivalence makes classically equivalent formulas have
the same conditional theory and normal worlds. Each original coordinate
can therefore be replaced by the chosen coordinate with identical rank,
label and satisfaction; every possible counterstate and the explicit
minimal predecessor remain. The restricted model is finite and preserves
the represented relation and smoothness. This is an actual finite-model
construction, not the false statement that finite vocabulary makes
syntactic formula sets finite.
\end{proof}
\end{annexitem}

\begin{annexitem}{annex:p08.rational.nonwellfounded}{thm:p08-rational-nonwellfounded}{Definable minima despite a descending chain}
\begin{proof}
Use the total order consisting of a bottom rank for \(s_\infty\)
and the reverse order of \(\omega\) above it. It induces exactly
the displayed preference. A finitary formula uses only finitely many
\(p_j\). For all sufficiently large finite \(i\), every such atom
is false at \(s_i\), exactly as at \(s_\infty\). Consequently,
if the formula has infinitely many finite-index states, it also has
\(s_\infty\), its minimum. If it omits \(s_\infty\), its set
is finite, and a nonempty such set has a maximum index, which is its
minimum in preference. Thus every nonempty definable set has a least
rank and the model is smooth and ranked.

The \(p_i\)-states are \(s_0,\ldots,s_i\), with unique minimum
\(s_i\). This label satisfies \(p_{i+1}\). The minimum of
\(p_{i+1}\) is \(s_{i+1}\), which falsifies \(p_i\).
The set for \(p_0\) is nonempty, so its assertion to \(\bot\)
fails. Now consider any smooth preferential model of this same relation.
There is a minimal \(p_0\)-state \(t_0\), and it satisfies
\(p_1\). It cannot be minimal for \(p_1\), because all those
minima falsify \(p_0\). Smoothness supplies \(t_1\prec t_0\)
minimal for \(p_1\). Repeating produces
\(t_{i+1}\prec t_i\), with \(t_i\) minimal for \(p_i\).
Under the usual choice background this dependent selection is a
descending chain. Equivalently, if the order were well founded,
choose a minimal point in the nonempty union of all \(p_i\)-sets;
the same argument gives a predecessor in that union, a contradiction.
Thus no well-founded representation exists.
\end{proof}
\end{annexitem}

\begin{annexitem}{annex:p08.rational.intersection}{thm:p08-rational-ranked-intersection}{Two ranked relations, their nonrational intersection and the selected closure}
\begin{proof}
Give each of two models two states at ranks 0 and 1. In coordinate
order \((p,q,a)\), the first has labels \((1,1,0)\) at rank 0
and \((1,0,1)\) at rank 1. The second has labels \((1,1,1)\)
at rank 0 and \((1,0,0)\) at rank 1. Both are finite smooth
ranked models, and their least \(p\)-state satisfies \(q\).
The first refutes \(p\wedge a\mathrel{|\!\sim}q\), since
its only \(p\wedge a\)-state is its rank-1 state.
The second refutes \(p\mathrel{|\!\sim}\neg a\), since
its least \(p\)-state satisfies \(a\).
Their relation intersection therefore has the two premises of Rational
Monotonicity, with antecedent \(p\), added information \(a\),
and consequent \(q\), but lacks the conclusion. Each constituent
relation is rational; their intersection is not.

For rational closure of the single-default base, its material
counterpart \(p\to q\) permits \(p\), so the default is
nonexceptional and \(C_1=\varnothing\). It also permits
\(p\wedge a\), giving this antecedent rank 0. It forbids
\(p\wedge a\wedge\neg q\), but the empty next base permits
that violation, giving it rank 1. The proved admissible selection
criterion therefore includes \(p\wedge a\mathrel{|\!\sim}q\).
The explicit first model shows it is not true in all ranked models
of the base. No general equality of ranked and preferential entailment
has been presumed in this separating proof.
\end{proof}
\end{annexitem}

\begin{annexitem}{annex:p08.rational.material}{thm:p08-rational-material-exception}{Materialization does exactly the top-level exception test}
\begin{proof}
If \(\widetilde K\not\models\delta\), choose a valuation
\(m\models\widetilde K\) falsifying \(\delta\).
The one-state model with label \(m\) and empty preference is smooth.
For each \(\alpha\mathrel{|\!\sim}\beta\in K\), either
\(m\) falsifies \(\alpha\), so the assertion is vacuous, or
material satisfaction makes it satisfy \(\beta\). Thus it models
\(K\) but refutes \(\top\mathrel{|\!\sim}\delta\).
Preferential soundness or semantic completeness therefore rules out
that assertion in \(K^p\).

Conversely, if \(\widetilde K\models\delta\), classical compactness
supplies finitely many members \(\alpha_j\to\beta_j\) jointly
entailing \(\delta\). Applying S to each original assertion with
outer antecedent \(\top\) gives
\(\top\mathrel{|\!\sim}\alpha_j\to\beta_j\).
Finite And and Right Weakening give
\(\top\mathrel{|\!\sim}\delta\).
If the finite subset is empty, \(\delta\) is a tautology and
Reflexivity plus Right Weakening gives the same result. Setting
\(\delta=\neg\alpha\) proves the exception test. The single-state
argument and compactness apply to arbitrary set-sized bases; the
algorithmic finiteness claims are not inferred from this general lemma.
\end{proof}
\end{annexitem}

\begin{annexitem}{annex:p08.rational.selection}{thm:p08-rational-source-selection}{The admissible criterion and both whole-relation preference conditions}
\begin{proof}
In this proof only, let \(h(\varphi)\) denote its ordinal exception
rank for \(K\), with \(\infty\) for no rank, and let \(T\)
be the relation defined by the theorem's criterion. The marker is above
all ordinal ranks that occur; the clause for \(h(\varphi)=\infty\)
is separate from strict comparison. The material exception lemma shows
that \(h(\varphi\vee\psi)=\min(h(\varphi),h(\psi))\),
that classically equivalent formulas have equal ranks, and that
\(h(\varphi\wedge\psi)\geq h(\varphi)\).
Indeed, at every stage a material theory entails the negation of a
disjunction iff it entails both negations, and entailment of
\(\neg\varphi\) implies entailment of
\(\neg(\varphi\wedge\psi)\). The bases decrease, so once a
formula is nonexceptional it remains so. These statements also cover
limit stages and no-rank formulas.

These identities prove rationality explicitly. Reflexivity holds because
\(\varphi\wedge\neg\varphi\) has no rank. Left Logical
Equivalence preserves both tested ranks. For Right Weakening, if
\(\models\beta\to\gamma\), the violation
\(\alpha\wedge\neg\gamma\) entails
\(\alpha\wedge\neg\beta\), so its rank is no smaller.
For And, the violation of the conjunction is the disjunction of the
two violations, both of rank above \(h(\alpha)\).
For Or, the violation is the disjunction of the two premise violations;
each is above its respective antecedent rank, so their minimum is
above the minimum of those antecedent ranks. If an antecedent has no
rank its violation also has no rank, and the same argument or the
separate vacuity clause applies.
If \(\alpha\mathrel{|\!\sim}_T\beta\) and
\(h(\alpha)<\infty\), the identity
\(\alpha=(\alpha\wedge\beta)\vee(\alpha\wedge\neg\beta)\)
makes \(h(\alpha\wedge\beta)=h(\alpha)\).
Together with the rank inequality for subsets, this proves Cautious
Monotonicity. If instead
\(\alpha\mathrel{|\!\not\sim}_T\neg\beta\),
the criterion says \(h(\alpha)<\infty\) and
\(h(\alpha\wedge\beta)=h(\alpha)\). The violation of any
\(\gamma\) already consequential from \(\alpha\) still has
higher rank after intersecting with \(\beta\), proving Rational
Monotonicity. The six native P rules and RM are therefore all satisfied.

Transfinite induction on survival in the decreasing bases gives:
an original assertion belongs to \(C_\sigma\) iff its antecedent
has rank at least \(\sigma\), or has no rank. At successors this
is the definition of \(E\); at limits it is the intersection
condition. If \(\alpha\mathrel{|\!\sim}\beta\in K\)
has antecedent rank \(\sigma\), it belongs to \(C_\sigma\).
Its material counterpart makes \(\alpha\wedge\neg\beta\)
exceptional there, hence of strictly higher rank or no rank. If its
antecedent has no rank, the vacuity clause applies. Thus \(K\subseteq T\).
By admissibility, all no-rank antecedents are already inconsistent
for \(K^p\); every rational extension of \(K\) contains every
assertion with such an antecedent.

Let \(R\ne T\) be a rational extension of \(K\). There is a least
ordinal \(\tau\) at which the relations differ on an antecedent
rank, since they agree on no-rank antecedents. They agree on all
assertions whose antecedents have rank below \(\tau\).
Put \(H=\{\varphi:h(\varphi)\geq\tau,\ h(\varphi)<\infty\}\).
For ranked formulas in \(T\), its source abbreviation
\(\varphi<\psi\) is exactly \(h(\varphi)<h(\psi)\):
the assertion is \(\varphi\vee\psi\mathrel{|\!\sim}_T\neg\psi\),
whose violation is \(\psi\). Represent \(R\) by the proved
general ranked representation. For a formula consistent in \(R\),
write \(\lambda(\varphi)\) for its least occupied layer in this
model; inconsistent formulas have no states. For consistent formulas,
the source abbreviation \(\varphi<\psi\) in \(R\) means
\(\lambda(\varphi)<\lambda(\psi)\).

First suppose every formula in \(H\) is inconsistent in \(R\).
Then \(R\) contains all assertions with ranked antecedents at least
\(\tau\), so \(T\setminus R\) is empty. Choose any antecedent
\(\gamma\) of rank \(\tau\) occurring in the difference.
The assertion \(\gamma\mathrel{|\!\sim}\bot\) is in
\(R\setminus T\). All \(T\)-antecedents strictly preferred to
\(\gamma\) have rank below \(\tau\), where the relations agree.
This witnesses condition 1; condition 2 is vacuous.

Next suppose the least layers of the \(R\)-consistent formulas in
\(H\) have a least member. Choose \(\alpha\in H\) attaining it,
and write \(a=\lambda(\alpha)\). To use a smooth tail model, replace
the ranked representation by its upward-copy expansion: at each
existing layer \(b\), include a copy of every world that occurred at
some layer at most \(b\). This changes no formula's least layer or
the labels at that layer. After cutting off all layers below \(a\),
a formula whose original least layer was below \(a\) has copies
at \(a\), and a formula whose least layer was at least \(a\)
keeps its original minima. The tail is therefore smooth and ranked.
No antecedent of a default in \(C_\tau\) occurs below \(a\):
those antecedents have rank at least \(\tau\) or no rank; the
latter have no states in \(R\), by admissibility. The tail consequently
still satisfies \(C_\tau\).
Every formula \(\beta\) with \(h(\beta)>\tau\) is exceptional
for \(C_\tau\), so all global minima of the tail falsify it.
Thus it has no states at layer \(a\), and if it has states its
least layer is strictly above \(a\). Applying this to \(\alpha\)
shows \(h(\alpha)=\tau\).

If \(\xi\mathrel{|\!\sim}\eta\in T\setminus R\), its
antecedent is \(R\)-consistent with rank at least \(\tau\), and
its failure gives a minimum \(\xi\)-state falsifying \(\eta\).
Hence \(\lambda(\xi)=\lambda(\xi\wedge\neg\eta)\).
Its violation has exception rank above \(\tau\), so the previous
paragraph gives \(\lambda(\xi)>a\).
There is an assertion in \(R\setminus T\) whose antecedent has
exception rank \(\tau\) and least \(R\)-layer \(a\): if the
first difference at \(\tau\) is an extra assertion with antecedent
\(\xi\) at layer \(a\), use it. If that antecedent instead has
greater least layer or is inconsistent in \(R\), use
\(\alpha\vee\xi\mathrel{|\!\sim}_R\neg\xi\).
It holds in \(R\), but fails in \(T\), since both formulas have
exception rank \(\tau\). If the difference at \(\tau\) is a
missing \(T\)-assertion, the same comparison works by the preceding
violation argument. Call this witness \(\gamma\mathrel{|\!\sim}\delta\).
Its lower \(T\)-antecedents have rank below \(\tau\), proving
condition 1. Its least \(R\)-layer is \(a\), strictly below
every antecedent of an assertion in \(T\setminus R\), so the
same witness proves condition 2 for every such assertion.

Finally suppose the consistent formulas in \(H\) have no least
\(R\)-layer. Choose \(\gamma\) of exception rank \(\tau\).
If it is consistent in \(R\), choose \(\delta\in H\) consistent
at a lower layer; if it is inconsistent, choose any consistent
\(\delta\in H\). In either case
\(\delta\vee\gamma\mathrel{|\!\sim}_R\neg\gamma\)
is in \(R\setminus T\), with antecedent exception rank \(\tau\).
Lower \(T\)-antecedents again lie below \(\tau\), proving
condition 1. For condition 2 fix any
\(\xi\mathrel{|\!\sim}\eta\in T\setminus R\).
It is \(R\)-consistent and belongs to \(H\). Among the
\(R\)-consistent formulas in \(H\) with layer at most
\(\lambda(\xi)\), choose \(\varepsilon\) of least exception
rank; this is possible because that nonempty set of ranks is a set
of ordinals. There is a consistent \(\zeta\in H\) at a lower
layer than \(\varepsilon\), since there is no least layer.
Minimality of the chosen exception rank gives
\(h(\zeta)\geq h(\varepsilon)\).
Thus \(\zeta\vee\varepsilon\mathrel{|\!\sim}_R\neg\varepsilon\)
lies in \(R\setminus T\): the least \(R\)-layer is that of
\(\zeta\), but in \(T\) the consequent fails at the
\(\varepsilon\) rank. Its antecedent layer is strictly below
\(\lambda(\varepsilon)\leq\lambda(\xi)\), so its antecedent
is strictly preferred to \(\xi\) in \(R\). This supplies
condition 2 for every missing assertion, including the case in which
there is no least layer. Both conditions therefore give \(T\prec R\).
Since \(R\) was an arbitrary different rational extension,
\(T\) satisfies the source definition of rational closure.
For uniqueness, the preference definition is asymmetric without needing
its full transitivity theorem. If \(T\prec R\) had a condition-1
witness \(\alpha\mathrel{|\!\sim}\beta\in R\setminus T\)
and also \(R\prec T\), condition 2 of the latter would supply
\(\rho\mathrel{|\!\sim}\eta\in T\setminus R\) with
\(\rho<\alpha\) in \(T\). Condition 1 of the former says
exactly that this assertion must belong to \(R\), a contradiction.
Two different closures would each be preferable to the other, so
there cannot be two.
\end{proof}
\end{annexitem}

\begin{annexitem}{annex:p08.rational.admissibility}{thm:p08-rational-finite-admissibility}{Stabilization, completely exceptional bases and admissibility}
\begin{proof}
The sequence decreases within a set of \(n\) assertions. Every
nonstable step deletes at least one assertion, so for some \(N\leq n\)
we have \(C_N=C_{N+1}\). All later stages equal \(C_N\), including
limit stages. An assertion belongs to \(C_i\) iff it originally
belongs to \(K\) and its antecedent has rank at least \(i\),
including no rank: prove this by induction, since survival at the
next stage is exactly exceptionality at the present one.
The exception tests decrease monotonically as the bases decrease.
Thus a formula with no earlier rank is no-rank exactly when it remains
exceptional for \(C_N\).

Let \(\Gamma\) be the disjunction of the finitely many antecedents
of \(C_N\), with \(\Gamma=\bot\) if there are none.
Since every one of these antecedents is exceptional for \(C_N\),
the material lemma gives
\(\widetilde C_N\models\neg\Gamma\).
In a smooth preferential model of \(C_N\), suppose a
\(\Gamma\)-state existed. Smoothness supplies a minimal one
\(s\). For every default whose antecedent is true at \(s\),
the state is minimal for that antecedent, since any predecessor
satisfying it would belong to \(\widehat\Gamma\). The default
therefore makes its consequent true at \(s\). Defaults with false
antecedents have true material implications. Hence the label of
\(s\) satisfies every member of \(\widetilde C_N\), and also
\(\Gamma\), contrary to the displayed classical entailment.
There are consequently no \(\Gamma\)-states. Every state's label
then satisfies all of \(\widetilde C_N\), by false antecedents.

A no-rank formula \(\alpha\) has
\(\widetilde C_N\models\neg\alpha\). No state in any model
of \(C_N\) can satisfy it. Every model of \(K\) is a model of
\(C_N\subseteq K\), so it satisfies
\(\alpha\mathrel{|\!\sim}\bot\).
The preferential completeness theorem puts this assertion in \(K^p\),
proving admissibility. Finiteness was used for a formula expressing the
union of all persistent antecedents, not to assume that all of
\(\mathcal U\), \(\mathcal L\), or a given model is finite.
\end{proof}
\end{annexitem}

\begin{annexitem}{annex:p08.rational.realization}{thm:p08-rational-finite-realization}{World ranks and every direction of the criterion}
\begin{proof}
The material theories \(\widetilde C_i\) decrease with \(i\),
so their model sets increase. At stage \(i\), a formula
\(\alpha\) is nonexceptional iff
\(\widetilde C_i\cup\{\alpha\}\) is satisfiable, by the
material lemma. If \(h(\alpha)=i<\infty\), such a witness
world \(m\) belongs to \(\mathcal U_*\), since
\(C_N\subseteq C_i\). It has \(k(m)\leq i\).
No \(\alpha\)-world has smaller \(k\), because that would
make \(\alpha\) nonexceptional at an earlier stage. Thus its
least state rank is exactly \(i\). If it has no rank, the stable
material theory entails its negation, and \(\mathcal U_*\) has
no such world. Conversely absence of any \(\alpha\)-world in
\(\mathcal U_*\) implies exceptionality at \(N\), hence no rank.

If \(\alpha\mathrel{|\!\sim}\beta\in K\) and
\(h(\alpha)=i<\infty\), the retention observation puts this
assertion in \(C_i\). All minimum \(\alpha\)-worlds satisfy
\(\widetilde C_i\) and \(\alpha\), hence \(\beta\).
If \(h(\alpha)=\infty\), the assertion is vacuous. Therefore
the model satisfies \(K\). The ranking maps into the finite
ordered set \(\{0,\ldots,N\}\); every nonempty formula set
has a least occupied layer, giving smoothness and ranked soundness.

For a finite-rank antecedent, its assertion holds exactly when no
state at its least layer satisfies \(\alpha\wedge\neg\beta\).
The latter set is a subset of the antecedent set, so it has either
strictly larger least rank or no states. By the first paragraph these
are exactly \(h(\alpha)<h(\alpha\wedge\neg\beta)\).
For a no-rank antecedent there are no states, so all assertions hold.
These cases prove both directions of the criterion. In particular,
\(\infty<\infty\) is not used: a no-rank antecedent is the
separate vacuity clause. Identification with the source's abstract
least-rational-extension definition uses the separately proved selection
theorem, not an inference from rationality of this one model.
\end{proof}
\end{annexitem}

\begin{annexitem}{annex:p08.rational.algorithm}{thm:p08-rational-finite-algorithm}{The stopping stage, the persistent case and the resource count}
\begin{proof}
As long as the current \(C_i\) is not stable and \(\alpha\)
is exceptional, the loop advances to \(C_{i+1}\).
If \(\alpha\) has rank \(i\), it reaches exactly that stage
and stops. At this stage its violation \(\alpha\wedge\neg\beta\)
is exceptional iff its rank is strictly greater than \(i\),
including no rank. This is precisely the finite-antecedent criterion.
If \(\alpha\) has no rank, the loop instead stops at \(C_N\).
This material theory entails \(\neg\alpha\), hence also
\(\neg(\alpha\wedge\neg\beta)\). The final test answers yes,
as required by the separate no-rank clause. These are all possible
stopping cases, because each nonstable iteration removes a default.

For every tested formula \(\delta\), exceptionality means
\(\widetilde C_i\models\neg\delta\), equivalently no valuation
satisfies the conjunction of its material implications and \(\delta\).
This equivalence also holds for empty \(C_i\).
At a stage at most \(n\) antecedents are tested to form \(E(C_i)\),
plus the query antecedent and the final violation test as needed.
There are at most \(n\) strict decreases, hence \(O(n^2)\) tests,
each with formula size bounded polynomially by the input sizes.
The procedure is thus deterministic with a polynomial number of SAT
oracle calls. No deterministic polynomial algorithm for unrestricted
SAT has been established here. If each tested conjunction is Horn,
the standard Horn satisfiability procedure can replace that oracle and
the same count gives deterministic polynomial time.
\end{proof}
\end{annexitem}

\begin{annexitem}{annex:p08.rational.minimum}{thm:p08-rational-pointwise-minimum}{Exactly possible labels and pointwise minimality}
\begin{proof}
Every label in a smooth model of \(K\) lies in \(\mathcal U_*\)
by the stable-base proof: persistent antecedents are absent, so the
label satisfies \(\widetilde C_N\). Conversely the constructed
model contains every valuation in \(\mathcal U_*\) and satisfies
\(K\). This proves the possible-label assertion.

On the finite vocabulary, each valuation has a characteristic
conjunction of literals. Its state set in a smooth ranked model has
a least rank whenever the valuation occurs. Retain one copy at that
rank for each valuation and delete all other copies. Formula truth
depends only on the label; every deleted copy has a retained copy
with the same label at no larger rank. Consequently each formula's
minimum rank and the labels at it are unchanged. There are finitely
many retained states, whose occupied layers can be numbered
\(0,1,\ldots,d\) in their original order. Call the resulting
valuation rank \(f(m)\). It preserves the model's relation.

We prove by induction on \(i\leq d\) that
\(f(m)\leq i\) implies \(k(m)\leq i\).
For \(i=0\), a world of global least rank satisfies every material
default of \(K\): if its antecedent is true there, it is a minimal
state for that antecedent. Thus \(k(m)=0\).
For the inductive step, consider any assertion
\(\gamma\mathrel{|\!\sim}\delta\in C_i\).
Its antecedent has exception rank at least \(i\), or no rank.
No \(\gamma\)-world can have \(f<i\), since the induction
hypotheses would give \(k<i\), contradicting the exact formula
rank realization. If a world with \(f(m)\leq i\) satisfies
\(\gamma\), it therefore has \(f(m)=i\) and is minimal for
\(\gamma\). The model of \(K\) makes it satisfy \(\delta\).
If it falsifies \(\gamma\), the material default is also true.
Hence it satisfies every material member of \(C_i\), and
\(k(m)\leq i\). For \(i>N\), the conclusion follows directly
from \(k(m)\leq N<i\), so all indices are covered.
Taking \(i=f(m)\) proves pointwise minimality.
Since the displayed \(k\) itself supplies a complete model, it is
the least member; any other least ranking must equal it at every
valuation. The proof compared complete normalized valuation models,
not arbitrary model cardinalities or the source's relation preference.
\end{proof}
\end{annexitem}

\begin{annexitem}{annex:p08.rational.penguin}{thm:p08-rational-penguin-example}{The full valuation calculation and inherited-property failure}
\begin{proof}
The initial material conjunction is
\((p\to b)\wedge(p\to\neg f)\wedge(b\to f)\),
which forces \(p=0\) and \(b\to f\). Thus \(p\) is
exceptional, but \(b\) is not, as witnessed by \((b,p,f)=(1,0,1)\).
Deleting just the nonexceptional bird default gives the displayed
\(C_1\). Its material theory permits \((1,1,0)\), so \(p\)
is now nonexceptional and both remaining assertions are removed:
\(C_2=\varnothing\). Every valuation is allowed in \(\mathcal U_*\).

In coordinate order \((b,p,f)\), all eight ranks are
\[
\begin{array}{c|cccccccc}
m&000&001&010&011&100&101&110&111\\\hline
k(m)&0&0&2&2&1&0&1&2
\end{array}
\]
The minimum bird rank is 0 at \(101\); the minimum penguin rank
is 1 at \(110\); the minimum flying-penguin rank is 2.
The minimum nonflying-bird rank is 1 at both \(100,110\), with
different penguin values. The asserted consequences and both failures
therefore follow directly from these minima.
Adding a fresh atom to the language alone leaves every material stage
unchanged. Each old world has extensions with either value at the same
rank, so \(p\wedge a\) still has minimum rank 1 and those minima
falsify \(f\). This justifies the particular irrelevant-premise
conclusion; it is not arbitrary monotonic strengthening.

Finally add the default \(b\mathrel{|\!\sim}a\) to the base.
At stage 0, materialization still forces \(\neg p\); worlds with
\(b=1,f=1,a=1,p=0\) show \(b\) nonexceptional. Both bird
defaults disappear at stage 1, leaving only
\(p\mathrel{|\!\sim}b,p\mathrel{|\!\sim}\neg f\).
The minimum penguin worlds then have \(b=1,p=1,f=0\) with
either value of \(a\). Hence neither value of \(a\) is a
conditional consequence of \(p\). This proves the specified
generic-property inheritance limitation by actual ranks.
\end{proof}
\end{annexitem}

\begin{annexitem}{annex:p08.rational.nixon}{thm:p08-rational-nixon-example}{A conflict at rank 1 is not a no-rank antecedent}
\begin{proof}
For the first base the material theory is
\((r\to\neg q)\wedge(s\to q)\).
It permits \(r\wedge\neg s\wedge\neg q\) and
\(s\wedge\neg r\wedge q\), so neither default is exceptional.
It forbids \(r\wedge s\), so this conjunction is exceptional
at stage 0. At stage 1 the base is empty, and every classically
satisfiable formula is nonexceptional. The conjunction and both
of its \(q\)-branches thus have rank 1. Neither branch is above
the antecedent's least layer, so neither conditional conclusion holds.
The fresh atom \(a\) imposes no material restriction; minimum
\(r\wedge a\)-states still have rank 0 and falsify \(q\).

For \(K'\), materialization is
\((p\to q)\wedge(p\to\neg q)\), equivalent to \(\neg p\).
Both antecedents are exceptional, so \(E(K')=K'\) and the base
is already stable. The possible worlds are exactly \(\neg p\)
worlds, all at rank 0. No \(p\)-state exists, so every assertion
with antecedent \(p\) holds, including both consequents and
\(\bot\). The first case had actual states at a higher finite
layer; the second excludes those states altogether.
\end{proof}
\end{annexitem}
